\Needspace{7\baselineskip}\Article{306}
The proof of the preceding theorem will be found entirely analogous to the proofs in Articles \(45,49\). Indeed, the theory of multiplication of classes has a very close affinity in every respect with the subject treated in Section III. But the limits of this work do not let us pursue this theory as fully as it deserves. So we will add only a few observations here, also leaving out proofs that would require a longer preparation, and reserve a fuller investigation for another occasion.

I.\ If the series \(K,C,2C,3C\), and so on, is continued beyond \((m-1)C\), the same classes appear again:
\[mC=K,\quad(m+1)C=C,\quad(m+2)C=2C\text{ etc.}\]
More generally, regarding \(K\) as \(0C\) for convenience, the classes \(gC,g'C\) are identical or different according as \(g\) and \(g'\) are congruent or incongruent modulo \(m\). Thus the class \(nC\) is always identical with the principal class \(K\).

II.\ We will call the collection of classes \(K,C,2C\ldots(m-1)C\), denoted above by \(\mathfrak C\), the \textit{period} of the class \(C\). This expression must not be confused with the \textit{periods of reduced forms} of positive nonsquare determinant treated in Article \(186\) and the following articles. Clearly, composing any number of classes in the same period produces a class also contained in that period:
\[gC+g'C+g''C\text{ etc.}=(g+g'+g''+\text{etc.})C\]


III.\ Since \(C+(m-1)C=K\), the classes \(C\) and \((m-1)C\) are opposite; so are \(2C\) and \((m-2)C\), \(3C\) and \((m-3)C\), and so on. If \(m\) is even, the class \(\tfrac12mC\) is therefore opposite to itself and so \textit{ambiguous}. Conversely, if \(\mathfrak C\) contains another ambiguous class besides \(K\), say \(gC\), then \(gC=(m-g)C\), giving \(g=m-g=\tfrac12m\). Thus, if \(m\) is even, there can be no ambiguous class other than \(K\) and \(\tfrac12mC\);
\SourceLeaf{377}
if \(m\) is odd, none other than \(K\) can be contained in \(\mathfrak C\).

IV.\ Suppose the period of a class \(hC\) contained in \(\mathfrak C\) is
\[K,hC,2hC,3hC\ldots(m'-1)hC\]
Then \(m'h\) is clearly the smallest multiple of \(h\) divisible by \(m\). So if \(m\) and \(h\) are coprime, \(m'=m\), and the two periods contain the same classes arranged in different orders. More generally, if \(\mu\) denotes the greatest common divisor of \(m,h\), then \(m'=\frac m\mu\). Thus the number of classes in the period of any class from \(\mathfrak C\) is either \(m\) or a proper divisor of \(m\). More precisely, as many classes in \(\mathfrak C\) have periods of \(m\) terms as there are numbers among \(0,1,2\ldots m-1\) coprime to \(m\), that is, \(\varphi m\), using the notation of Article \(39\). In general, as many classes in \(\mathfrak C\) have periods of \(\frac m\mu\) terms as there are numbers among \(0,1,2\ldots m-1\) having \(\mu\) as their greatest common divisor with \(m\). Their number is readily seen to be \(\varphi\frac m\mu\). So if \(m=n\), that is, if \(\mathfrak C\) contains the \textit{whole} principal genus, this genus contains exactly \(\varphi n\) classes whose periods include the whole genus, and \(\varphi e\) classes whose periods have \(e\) terms, where \(e\) is any divisor of \(n\). This conclusion holds in general whenever the principal genus contains a class whose period has \(n\) terms.

V.\ Under the same assumption, the system of classes of the principal genus is most conveniently arranged by adopting a class whose period has \(n\) terms as a kind of \textit{base}, and placing the classes of the principal genus in the order in which they occur in that period. If we assign the \textit{index} \(0\) to the principal class, \(1\) to the class chosen as base, and so on, then merely adding indices tells us which class results from composing any classes of the principal genus. Here is an example for the determinant \(-356\), with the class \((9,2,40)\) chosen as base:
\begin{center}\begin{tabular}{rl|rl|rl}
\(0\) & \((1,0,356)\) & \(4\) & \((20,8,21)\) & \(8\) & \((20,-8,21)\)\\[2pt]
\(1\) & \((9,2,40)\) & \(5\) & \((17,1,21)\) & \(9\) & \((8,2,45)\)\\[2pt]
\(2\) & \((5,2,72)\) & \(6\) & \((4,0,89)\) & \(10\) & \((5,-2,72)\)\\[2pt]
\(3\) & \((8,-2,45)\) & \(7\) & \((17,-1,21)\) & \(11\) & \((9,-2,40)\)\\[2pt]
\end{tabular}\end{center}

\SourceLeaf{378}

VI.\ The analogy with Section III, together with examination of more than \(200\) negative determinants and far more positive nonsquare ones, may seem to make it highly probable that this assumption holds for \textit{all} determinants. Yet that conclusion would be false, and extending the table of classifications would refute it. For brevity, let us call determinants for which the whole principal genus can be included in a single period \textit{regular}, and the others, for which this is impossible, \textit{irregular}. This subject seems to belong to the deepest mysteries of higher arithmetic and to leave room for extremely difficult investigations. We can illustrate it here only with a few observations, preceded by the following general one.

VII.\ If the principal genus contains classes \(C,C'\) whose periods consist of \(m,m'\) classes, and \(M\) is the smallest number divisible by \(m\) and \(m'\), then the same genus also contains classes whose periods have \(M\) terms. Resolve \(M\) into two coprime factors \(r,r'\), the former \((r)\) dividing \(m\) and the latter \((r')\) dividing \(m'\), as in Article \(73\). Then the class \(\frac mrC+\frac{m'}{r'}C'=C''\) has the required property. For suppose the period of \(C''\) has \(g\) terms. Then
\[\begin{aligned}K=grC''&=gmC+\frac{grm'}{r'}C'\\&=K+\frac{grm'}{r'}C'=\frac{grm'}{r'}C'\end{aligned}\]
so \(\frac{grm'}{r'}\) must be divisible by \(m'\), or \(gr\) by \(r'\), and therefore also \(g\) by \(r'\). In exactly the same way \(g\) is found divisible by \(r\), and so also by \(rr'=M\). But clearly \(MC''=K\), so \(M\) is also divisible by \(g\). Thus necessarily \(M=g\). It readily follows that the \textit{largest} number of classes contained in any period for a given determinant is divisible by the number of classes in every other period of a class from the same principal genus. This also gives a method of finding a class whose period is as large as possible, and so, for a regular determinant, includes the whole principal genus. It is entirely analogous to the method of Articles \(73,74\), though in practice various devices can shorten the work. The quotient of \(n\) by the number of classes in the largest period is \(1\) for regular determinants, but for irregular ones is always an integer greater than \(1\). For the latter it is particularly convenient for expressing the different kinds of irregularity, and may therefore be called the \textit{irregularity exponent}.
\SourceLeaf{379}

VIII.\ So far there is no general rule for distinguishing regular from irregular determinants in advance, especially since the latter include both primes and composite numbers. A few particular observations must therefore suffice. If the principal genus contains more than two ambiguous classes, the determinant is certainly irregular and its irregularity exponent even. If that genus contains just one or two, the determinant is either regular or at least has an odd irregularity exponent. All negative determinants of the form \(-(216k+27)\), except \(-27\) alone, are irregular, with irregularity exponent divisible by \(3\). The same holds for negative determinants of the forms \(-(1000k+75)\) and \(-(1000k+675)\), except \(-75\) alone, and for infinitely many others. If the irregularity exponent is a prime \(p\), or is at least divisible by \(p\), then \(n\) is divisible by \(pp\). Thus if \(n\) has no square divisor, the determinant is certainly regular. Only for positive \textit{square} determinants \(ee\) can we always tell in advance whether they are regular or irregular: regular when \(e\) is \(1\) or \(2\), an odd prime or a power of an odd prime, and irregular in all other cases. Among negative determinants, irregular ones grow steadily more frequent as the determinants grow. \textit{For example}, in the whole first thousand there are thirteen irregular ones: omitting the negative sign, \(576,580,820,884,900\), whose irregularity exponent is \(2\), and \(243,307,339,459,675,755,891,974\), whose irregularity exponent is \(3\). In the second thousand we have found \(13\) with exponent \(2\) and \(15\) with exponent \(3\); in the tenth thousand, \(31\) with exponent \(2\) and \(32\) with exponent \(3\). We cannot yet decide whether determinants with irregularity exponent greater than \(3\) occur below \(-10000\). Beyond this limit any given exponent can occur. It is highly probable that, as the determinants grow, the frequency of irregular negative determinants relative to regular ones approaches a constant ratio ever more closely. Determining that ratio would be very worthy of mathematicians' ingenuity. — For positive nonsquare determinants, irregular ones are much rarer. Certainly infinitely many have an even irregularity exponent, \textit{for example} \(3026\), for which it is \(2\). There also seems no doubt that ones with odd irregularity exponent exist, although we must admit that so far none has presented itself to us.

IX.\ For brevity we cannot discuss here the most convenient arrangement of the system of classes in the principal genus for an irregular determinant. We only observe that, since a single base does not suffice, two
\SourceLeaf{380}
or even more classes must be chosen, from whose multiplication and composition all the classes can be produced. This gives \textit{double or multiple indices}, serving much the same purpose as single indices for regular determinants. We will treat this more fully another time.

X.\ Finally, all the properties considered in this and the preceding article depend chiefly on the number \(n\), which is analogous to \(p-1\) in Section III. So this number deserves the closest attention, and it is most desirable to discover a general connection between it and its determinant. We think there is all the less reason to lose hope about this very important question, since we have already succeeded in expressing by an analytic formula the average value of the product of \(n\) and the number of genera, which can be determined in advance, at least for negative determinants, in Article \(302\).

\Needspace{7\baselineskip}\Article{307}
The investigations of the preceding articles cover only the classes of the principal genus. So they suffice both for positive determinants with just one genus altogether and for negative determinants with just one positive genus, if we choose not to consider the negative genus. It remains to add some observations on the other properly primitive genera as well.

I.\ If a genus \(G'\), different from the principal genus \(G\) of the same determinant, contains any ambiguous class, it contains as many as \(G\). Let the ambiguous classes in \(G\) be \(L\), \(M,N\), and so on, including the principal class \(K\), and those in \(G'\) be \(L',M',N'\), and so on. Denote the collection of the former by \(A\), and of the latter by \(A'\). Clearly all the classes \(L+L',\ M+L',\ N+L'\), and so on, are ambiguous and different, and belong to \(G'\), so must be contained in \(A'\). Thus the number of classes in \(A'\) certainly cannot be smaller than in \(A\). Likewise the classes \(L'+L'\), \(M'+L',\ N'+L'\), and so on, are different and ambiguous, belong to \(G\), and so are contained in \(A\). Thus the number of classes in \(A\) cannot be smaller than in \(A'\). The numbers of classes in \(A\) and \(A'\) must therefore be equal.

II.\ Since the total number of ambiguous classes equals the number of genera, by Articles \(261,287\) III, if \(G\) contains just one ambiguous
\SourceLeaf{381}
class, then clearly \textit{every} genus must contain one ambiguous class. If \(G\) contains two, half of all genera contain two each and the others none. Finally, if \(G\) contains more, say \(a\)\HistoricalNote{\ensuremath{\ast})}{This can happen only for irregular determinants, and \(a\) will always be a power of two.}, one \(a\)-th of all the genera contain \(a\) ambiguous classes each, and the others none.

III.\ In the case where \(G\) contains two ambiguous classes, let \(G,G',G''\), and so on, be the genera containing two each, and \(H,H',H''\), and so on, those containing none. Denote the former collection by \(\mathfrak G\) and the latter by \(\mathfrak H\). Composing two ambiguous classes always gives an ambiguous class, by Article \(249\). So composing two genera from \(\mathfrak G\) clearly always gives a genus from \(\mathfrak G\). It follows further that composing a genus from \(\mathfrak G\) with one from \(\mathfrak H\) gives a genus from \(\mathfrak H\). For if, \textit{for example}, \(G'+H\) belonged not to \(\mathfrak H\) but to \(\mathfrak G\), then \(G'+H+G'\) would also have to belong to \(\mathfrak G\). \textit{This is a contradiction}, since \(G'+G'=G\), and therefore \(G'+H+G'=H\). Finally, the genera \(G+H\), \(G'+H,\ G''+H\), and so on, together with \(H+H,\ H'+H,\ H''+H\), and so on, are clearly all different, and so coincide with \(\mathfrak G\) and \(\mathfrak H\) taken together. But by what we have just proved the genera \(G+H,\ G'+H,\ G''+H\), and so on, all belong to \(\mathfrak H\), and therefore exhaust this collection. So the remaining genera \(H+H,\ H'+H\), \(H''+H\), and so on, must all belong to \(\mathfrak G\). \textit{That is}, composing two genera from \(\mathfrak H\) always gives a genus from \(\mathfrak G\).

IV.\ If \(E\) is a class of a genus \(V\) different from the principal genus \(G\), then clearly \(2E,4E\), \(6E\), and so on, all belong to \(G\), while \(3E,5E,7E\), and so on, belong to \(V\). If the period of \(2E\) has \(m\) terms, then in the series \(E,2E,3E\), and so on, the class \(2mE\), and no earlier one, is clearly identical with \(K\). So the period of \(E\) has \(2m\) terms. Thus the number of terms in the period of any class outside the principal genus is either \(2n\) or a proper divisor of \(2n\), where \(n\) is the number of classes in each genus.

V.\ Let \(C\) be a given class of the principal genus \(G\), and \(E\) a class of a genus \(V\) whose doubling gives \(C\), which always exists by Article \(286\). Let all the ambiguous classes, properly primitive and of the same determinant, be \(K,K',K''\), and so on. Then \textit{all} classes whose
\SourceLeaf{382}
doubling gives \(C\) are \(E\ (=E+K),\ E+K',\ E+K''\), and so on. Denote their collection by \(\mathfrak Q\). Their number equals the number of ambiguous classes, or the number of genera. Clearly, as many classes in \(\mathfrak Q\) belong to the genus \(V\) as there are ambiguous classes in \(G\). Denoting this number by \(a\), we see that each genus contains either \(a\) classes from \(\mathfrak Q\) or none. Thus, when \(a=1\), each genus contains one class from \(\mathfrak Q\); when \(a=2\), half of all genera contain two classes each from \(\mathfrak Q\), and the rest none. The former half either coincides entirely with \(\mathfrak G\), in the same sense as in III above, and the latter with \(\mathfrak H\), or the latter coincides with \(\mathfrak G\) and the former with \(\mathfrak H\). — When \(a\) is larger still, one \(a\)-th of all the genera always contain classes from \(\mathfrak Q\), each containing \(a\) classes.

VI.\ Now suppose \(C\) is a class whose period has \(n\) terms. In the case where \(a=2\), and therefore \(n\) is even, it is easy to see that none of \(\mathfrak Q\) can belong to \(G\). For such a class would then be contained in the period of \(C\). If it were \(=rC\), or \(2rC=C\), then \(2r\equiv1\; (\mathrm{mod}.\ n)\), \textit{which is a contradiction}. Since \(G\) belongs to \(\mathfrak G\), all the classes \(\mathfrak Q\) must therefore be distributed among the genera \(\mathfrak H\). Now, for a regular determinant, \(G\) contains exactly \(\varphi n\) classes whose periods have \(n\) terms. So in the case \(a=2\), each genus \(\mathfrak H\) contains exactly \(2\varphi n\) classes whose periods include \(2n\) terms, and therefore both their own genus and the principal genus. When \(a=1\), each genus other than the principal one contains \(\varphi n\) such classes.

VII.\ On these observations we base the following method of arranging the system of \textit{all} properly primitive classes as conveniently as possible for any given regular determinant; irregular ones are set aside entirely. Choose any class \(E\) whose period contains \(2n\) terms, and therefore both its own genus, say \(V\), and the principal genus \(G\). Arrange the classes of these two genera in their order in that period. This already finishes the task when there are no more than these two genera, or when it seems unnecessary to add the others, \textit{for example} for a negative determinant with just two positive genera. If four or more genera are to be arranged, treat the others as follows. Let \(V'\) be one of the others and \(V+V'=V''\). Then \(V'\) and \(V''\) together contain two ambiguous classes, either one in each or two in one and none in the other. From these
\SourceLeaf{383}
choose any one \(A\). Composing \(A\) with each class in \(G\) and \(V\) clearly gives \(2n\) different classes belonging to \(V'\) and \(V''\), and therefore exhausting these genera. So these too can be arranged. — If genera remain beyond these four, let \(V'''\) be one of them, and \(V'''',V''''',V''''''\) the genera obtained by composing \(V'''\) with \(V,V'\) and \(V''\). These four genera \(V'''\ldots V''''''\) contain four ambiguous classes. Choosing one of them, \(A'\), and composing it with each class in \(G,V,V',V''\) clearly produces all classes in \(V'''\ldots V''''''\). — If more genera remain, continue in the same way until all are exhausted. Clearly, if the number of genera to be arranged is \(2^\mu\), we need exactly \(\mu-1\) ambiguous classes. Any class in these genera can be produced either by multiplication of \(E\), or by composing a class obtained from such multiplication with one or more ambiguous classes. Here are two examples illustrating these rules. We cannot add more here on the use of such an arrangement or the devices that can lighten the work.

\Needspace{12\baselineskip}\begin{center}I. Determinant \(-161\).\par\smallskip Four positive genera; four classes in each.\end{center}
\Needspace{9\baselineskip}
\begin{center}
\begin{minipage}[t]{0.47\linewidth}\centering
\(G\)\par\smallskip
{\small \(1,4;\ R7;\ R23\)}\par\smallskip
\begin{tabular}{@{}l@{}}
\((1,0,161)=K\)\\\((9,1,18)=2E\)\\\((2,1,81)=4E\)\\\((9,-1,18)=6E\)
\end{tabular}
\end{minipage}\hfill\begin{minipage}[t]{0.47\linewidth}\centering
\(V\)\par\smallskip
{\small \(3,4;\ N7;\ R23\)}\par\smallskip
\begin{tabular}{@{}l@{}}
\((3,1,54)=E\)\\\((6,-1,27)=3E\)\\\((6,1,27)=5E\)\\\((3,-1,54)=7E\)
\end{tabular}
\end{minipage}
\end{center}
\Needspace{9\baselineskip}
\begin{center}
\begin{minipage}[t]{0.47\linewidth}\centering
\(V'\)\par\smallskip
{\small \(3,4;\ R7;\ N23\)}\par\smallskip
\begin{tabular}{@{}l@{}}
\((7,0,23)=A\)\\\((11,-2,15)=A+2E\)\\\((14,7,15)=A+4E\)\\\((11,2,15)=A+6E\)
\end{tabular}
\end{minipage}\hfill\begin{minipage}[t]{0.47\linewidth}\centering
\(V''\)\par\smallskip
{\small \(1,4;\ N7;\ N23\)}\par\smallskip
\begin{tabular}{@{}l@{}}
\((10,3,17)=A+E\)\\\((5,2,33)=A+3E\)\\\((5,-2,33)=A+5E\)\\\((10,-3,17)=A+7E\)
\end{tabular}
\end{minipage}
\end{center}
\SourceLeaf{384}
\Needspace{12\baselineskip}\begin{center}II. Determinant \(-546\).\par\smallskip Eight positive genera; three classes in each.\end{center}
\Needspace{9\baselineskip}
\begin{center}
\begin{minipage}[t]{0.47\linewidth}\centering
\(G\)\par\smallskip
{\small \(1\text{ and }3, 8; R3; R7; R13\)}\par\smallskip
\begin{tabular}{@{}l@{}}
\((1,0,546)=K\)\\\((22,-2,25)=2E\)\\\((22,2,25)=4E\)
\end{tabular}
\end{minipage}\hfill\begin{minipage}[t]{0.47\linewidth}\centering
\(V\)\par\smallskip
{\small \(5\text{ and }7, 8; N3; N7; N13\)}\par\smallskip
\begin{tabular}{@{}l@{}}
\((5,2,110)=E\)\\\((21,0,26)=3E\)\\\((5,-2,110)=5E\)
\end{tabular}
\end{minipage}
\end{center}
\Needspace{9\baselineskip}
\begin{center}
\begin{minipage}[t]{0.47\linewidth}\centering
\(V'\)\par\smallskip
{\small \(1\text{ and }3, 8; N3; R7; N13\)}\par\smallskip
\begin{tabular}{@{}l@{}}
\((2,0,273)=A\)\\\((11,-2,50)=A+2E\)\\\((11,2,50)=A+4E\)
\end{tabular}
\end{minipage}\hfill\begin{minipage}[t]{0.47\linewidth}\centering
\(V''\)\par\smallskip
{\small \(5\text{ and }7, 8; R3; N7; R13\)}\par\smallskip
\begin{tabular}{@{}l@{}}
\((10,2,55)=A+E\)\\\((13,0,42)=A+3E\)\\\((10,-2,55)=A+5E\)
\end{tabular}
\end{minipage}
\end{center}
\Needspace{9\baselineskip}
\begin{center}
\begin{minipage}[t]{0.47\linewidth}\centering
\(V'''\)\par\smallskip
{\small \(1\text{ and }3, 8; N3; N7; R13\)}\par\smallskip
\begin{tabular}{@{}l@{}}
\((3,0,182)=A'\)\\\((17,7,35)=A'+2E\)\\\((17,-7,35)=A'+4E\)
\end{tabular}
\end{minipage}\hfill\begin{minipage}[t]{0.47\linewidth}\centering
\(V''''\)\par\smallskip
{\small \(5\text{ and }7, 8; R3; R7; N13\)}\par\smallskip
\begin{tabular}{@{}l@{}}
\((15,-3,37)=A'+E\)\\\((7,0,78)=A'+3E\)\\\((15,3,37)=A'+5E\)
\end{tabular}
\end{minipage}
\end{center}
\Needspace{9\baselineskip}
\begin{center}
\begin{minipage}[t]{0.47\linewidth}\centering
\(V'''''\)\par\smallskip
{\small \(1\text{ and }3, 8; R3; N7; N13\)}\par\smallskip
\begin{tabular}{@{}l@{}}
\((6,0,91)=A+A'\)\\\((19,9,33)=A+A'+2E\)\\\((19,-9,33)=A+A'+4E\)
\end{tabular}
\end{minipage}\hfill\begin{minipage}[t]{0.47\linewidth}\centering
\(V''''''\)\par\smallskip
{\small \(5\text{ and }7, 8; N3; R7; R13\)}\par\smallskip
\begin{tabular}{@{}l@{}}
\((23,11,29)=A+A'+E\)\\\((14,0,39)=A+A'+3E\)\\\((23,-11,29)=A+A'+5E\)
\end{tabular}
\end{minipage}
\end{center}\BookRule{0.2\linewidth}
